\documentclass[12pt]{article}
\usepackage[T1]{fontenc}
\usepackage[utf8]{inputenc}
\usepackage{lmodern}
\usepackage{textcomp}
\usepackage{enumitem}
\usepackage{geometry}
\geometry{margin=1in}
\begin{document}
\begin{center}
Answer Key - Chemistry - Undergraduate Level Assessment 5 \
\small (Answers are ordered exactly as in the question paper.)
\end{center}

\section*{Multiple Choice}
\begin{enumerate}[label=\arabic*.]
\item \textbf{Answer:} D
\\ \textit{Options:} A) Momentum is known exactly, but no information about position can be known.; B) Position is known exactly, but no information about momentum can be known.; C) No information about either position or momentum can be known.; D) Neither position nor momentum can be known exactly.
\\ \textit{Note:} The Heisenberg uncertainty principle states that the more precisely the position of a particle is known, the less precisely its momentum can be known, and vice versa, which applies to a quantum particle in a one-dimensional box, indicating that neither property can be determined exactly.
\item \textbf{Answer:} D
\\ \textit{Options:} A) Both 35 Cl and 37 Cl have I = 0.; B) The hydrogen atom undergoes rapid intermolecular exchange.; C) The molecule is not rigid.; D) Both 35 Cl and 37 Cl have electric quadrupole moments.
\\ \textit{Note:} The presence of electric quadrupole moments in both isotopes of chlorine (35 Cl and 37 Cl) leads to a lack of nuclear spin coupling with the hydrogen atoms in CHCl 3, resulting in the observed singlet in the 1 H NMR spectrum.
\item \textbf{Answer:} C
\\ \textit{Options:} A) Protons and neutrons have orbital and spin angular momentum.; B) Protons have orbital and spin angular momentum, neutrons have spin angular momentum.; C) Protons and neutrons possess orbital angular momentum only.; D) Protons and neutrons possess spin angular momentum only.
\\ \textit{Note:} Protons and neutrons possess intrinsic spin, which contributes to their overall angular momentum, while electrons are responsible for orbital angular momentum due to their movement in atomic orbitals.
\item \textbf{Answer:} B
\\ \textit{Options:} A) HCl / NaOH; B) H$_{3}$O+ / H$_{2}$O; C) O$_{2}$ / H$_{2}$O; D) H+ / Cl-
\\ \textit{Note:} H$_{3}$O+ and H$_{2}$O are classified as a conjugate acid-base pair because H$_{3}$O+ is the protonated form (conjugate acid) of H$_{2}$O, which can donate a proton to become H$_{3}$O+.
\item \textbf{Answer:} D
\\ \textit{Options:} A) The most common oxidation state for the lanthanide elements is +3.; B) Lanthanide complexes often have high coordination numbers (\textgreater{} 6).; C) All of the lanthanide elements react with aqueous acid to liberate hydrogen.; D) The atomic radii of the lanthanide elements increase across the period from La to Lu.
\\ \textit{Note:} The statement is incorrect because the atomic radii of the lanthanide elements actually decrease across the period from lanthanum (La) to lutetium (Lu) due to the increasing effective nuclear charge, which pulls the electrons closer to the nucleus.
\end{enumerate}

\section*{True / False}
\begin{enumerate}[label=\arabic*.]
\setcounter{enumi}{5}
\item \textbf{Answer:} False
\\ \textit{Options:} A) True; B) False
\\ \textit{Note:} The statement is false because CH 3 F and CH 3 Cl have different electronegativities and molecular environments, leading to distinct chemical shifts in their 1 H NMR spectra despite being measured in the same magnetic field.
\item \textbf{Answer:} False
\\ \textit{Options:} A) True; B) False
\\ \textit{Note:} The statement is false because the polarization of a material is not directly determined by a magnetic field strength of 100 T at standard conditions, as typical magnetic fields in these contexts are much lower, and the provided polarization value is not consistent with the expected behavior of materials under such extreme magnetic conditions.
\item \textbf{Answer:} True
\\ \textit{Options:} A) True; B) False
\\ \textit{Note:} Unpaired electrons in materials create a net magnetic moment that contributes to paramagnetism and can align in parallel to produce ferromagnetism, making their presence essential for both phenomena.
\item \textbf{Answer:} False
\\ \textit{Options:} A) True; B) False
\\ \textit{Note:} Lead is toxic to aquatic organisms and can accumulate in the food chain, leading to harmful effects on both ecosystems and human health, which makes it unsafe in aquatic environments.
\item \textbf{Answer:} False
\\ \textit{Options:} A) True; B) False
\\ \textit{Note:} Ammonia, NH 3, is not the strongest base in liquid ammonia because stronger bases, such as alkali metal amides, can exist in this solvent and can deprotonate NH 3 more effectively.
\end{enumerate}

\section*{Long Form Response}
\begin{enumerate}[label=\arabic*.]
\setcounter{enumi}{10}
\item \textbf{Answer:} The orbital angular momentum quantum number, l, plays a crucial role in determining the ease of electron removal during the ionization of an atom, such as aluminum. In aluminum, which has the electron configuration [Ne]3 $s^23$ p¹, the outermost electron resides in the 3 p orbital, characterized by l = 1. Electrons in higher l orbitals (like p or d) generally have higher energy and are less tightly bound compared to those in lower l orbitals (like s). Therefore, during ionization, the electron in the 3 p orbital is the most easily removed due to its higher energy state and the influence of its angular momentum, which results in increased shielding and reduced effective nuclear charge experienced by that electron. This understanding of the significance of l highlights how the shape and energy of electron orbitals directly affect ionization processes in atoms.
\\ \textit{Note:} The answer correctly identifies that the orbital angular momentum quantum number, l, indicates the shape and energy of the electron's orbital, with higher l values corresponding to higher energy and less tightly bound electrons, making the 3 p electron in aluminum the most easily removed during ionization.
\item \textbf{Answer:} The difference in first ionization energy between argon and neon primarily arises from their atomic structure and electron configuration. Neon, with its atomic number of 10, has a complete outer shell of eight electrons (1 $s^2$ 2 $s^2$ 2 p⁶), resulting in a strong effective nuclear charge that holds the electrons tightly, thereby requiring more energy to remove one. In contrast, argon, with an atomic number of 18, has a similar electron configuration but is larger in size, leading to increased electron shielding and a lower effective nuclear charge experienced by its outermost electrons. This reduced attraction between the nucleus and the outer electrons in argon means that less energy is needed to remove an electron, resulting in a lower first ionization energy compared to neon.
\\ \textit{Note:} The answer correctly identifies that neon's complete outer electron shell and smaller atomic radius result in a higher effective nuclear charge, making it more difficult to remove an electron compared to argon, which, despite having a similar electron configuration, has a larger atomic size that reduces the effective nuclear charge felt by its outer electrons.
\end{enumerate}

\vfill
\noindent\textit{End of Answer Key}
\end{document}